% Defense Presentation — José Gonçalves Lima Neto
% IME-USP, Computer Science — Defense: 11 June 2026
%
% Rebuilt from scratch for the defense, using the deposited thesis
% (archive/thesis-deposit-original-20260511.pdf) as the authoritative
% source and the IME defense template (Template.pdf) as the structure guide.
% Flow: title -> big picture/context -> research questions -> method ->
%        advances (usage, uncertainties/patterns, refinements) ->
%        implications -> future work/closing.
%
% Built on the IME-USP Beamer template (ccsl-usp/modelo-latex).

%%%%%%%%%%%%%%%%%%%%%%%%%%%%%%%%%%%%%%%%%%%%%%%%%%%%%%%%%%%%%%%%%%%%%%%%%%%%%%%%
%%%%%%%%%%%%%%%%%%%%%%%%%%% DOCUMENT CLASS & PACKAGES %%%%%%%%%%%%%%%%%%%%%%%%%%
%%%%%%%%%%%%%%%%%%%%%%%%%%%%%%%%%%%%%%%%%%%%%%%%%%%%%%%%%%%%%%%%%%%%%%%%%%%%%%%%

\documentclass[
  xcolor={hyperref,svgnames,x11names,table},
  hyperref={pdfencoding=unicode,plainpages=false,pdfpagelabels=true,breaklinks=true},
  brazilian,english,12pt,aspectratio=169,
]{beamer}

\usepackage{imegoodies}
\usepackage[slides,hidelinks]{imelooks}

\usetikzlibrary{positioning,arrows.meta,fit,backgrounds,calc,shadows.blur}

\graphicspath{{figures/},{template/logos/}}

\babeltags{br = brazilian, en = english}

% Roomy frames, large readable type: no navigation clutter.
\setbeamertemplate{navigation symbols}{}

% The full thesis title is long (must match the deposit); shrink the title
% font so it fits the 16:9 title slide without clipping above the frame.
\setbeamerfont{title}{size=\large,series=\bfseries}

%%%%%%%%%%%%%%%%%%%%%%%%%%%%%%%%%%%%%%%%%%%%%%%%%%%%%%%%%%%%%%%%%%%%%%%%%%%%%%%%
%%%%%%%%%%%%%%%%%%%%%%%%%% METADATA & TITLE SLIDE %%%%%%%%%%%%%%%%%%%%%%%%%%%%%%
%%%%%%%%%%%%%%%%%%%%%%%%%%%%%%%%%%%%%%%%%%%%%%%%%%%%%%%%%%%%%%%%%%%%%%%%%%%%%%%%

% Bibliography source (same keys as the deposited thesis) for inline citations.
\addbibresource{references.bib}

\title[Hypo-Stage]{Hypo-Stage: Tool-Supported Architectural Hypothesis
Engineering in Professional Software Development Teams --- A Case Study}

\author[Lima Neto]{José Gonçalves Lima Neto}

\institute{\textbf{Master's Dissertation Defense}\\
           Advisor: Prof.\ Dr.\ Paulo Roberto Miranda Meirelles\\
           Graduate Program in Computer Science\\
           IME -- University of São Paulo}

\date{11 June 2026}

% Title page background image and footer logos (IME template assets).
\titlebgimage{\includegraphics[width=\paperwidth,height=\paperheight]{bg-ime}}

\logos{%
  \includegraphics[height=.9\baselineskip]{ime-logo}%
  \quad\quad
  \raisebox{.15\baselineskip}{\includegraphics[height=.6\baselineskip]{usp-logo}}%
  \par
  \hspace{.87\paperwidth}\includegraphics[width=.05\paperwidth]{cc-by}\par
}

\qrcodeurl{https://github.com/ArchHypo/hypo-stage-masters-thesis}


%%%%%%%%%%%%%%%%%%%%%%%%%%%%%%%%%%%%%%%%%%%%%%%%%%%%%%%%%%%%%%%%%%%%%%%%%%%%%%%%
%%%%%%%%%%%%%%%%%%%%%%%%%%%%%%%%%% SLIDES %%%%%%%%%%%%%%%%%%%%%%%%%%%%%%%%%%%%%%
%%%%%%%%%%%%%%%%%%%%%%%%%%%%%%%%%%%%%%%%%%%%%%%%%%%%%%%%%%%%%%%%%%%%%%%%%%%%%%%%

\begin{document}

% ============================================================================
% Slide 1 — Title
% ============================================================================
% SCRIPT: "This work asks how professional software teams manage architectural
%          uncertainty when they are given real tool support --- and what
%          happens once they actually put it to use."
\maketitle
\note{%
  Script: This dissertation studies how teams can manage architectural
  uncertainty with tool support in a professional software development setting.\par
  \medskip
  Speaker notes: Greet the committee, state the title, your name, and the
  advisor. Frame the talk in one breath: the work is an empirical case study
  of Hypo-Stage --- a Backstage plugin for the ArchHypo technique --- used by
  real teams in industry. Keep this to under a minute and move quickly to the
  big-picture slide, which carries the rest of the framing.%
}


% ============================================================================
% Slide 2 — Context / big-picture overview  (CENTRAL VISUAL)
% ============================================================================
% SCRIPT: "Here is the whole talk in one picture: architectural uncertainty is
%          unavoidable, ArchHypo makes it explicit as falsifiable hypotheses,
%          and Hypo-Stage operationalizes that inside the developer portal ---
%          this case study asks how teams actually use it."
\begin{frame}{Context}
  \begin{center}
    \begin{tikzpicture}
      \node[rounded corners,fill=white,draw=black!25,line width=0.5pt,inner sep=8pt,
            blur shadow={shadow blur steps=6,shadow xshift=0.7pt,shadow yshift=-0.9pt}]
        {\includegraphics[height=0.85\textheight]{dissertation-big-picture}};
    \end{tikzpicture}
  \end{center}
\end{frame}
\note{%
  Script: The core idea is simple --- uncertainty exists, ArchHypo structures
  it, Hypo-Stage operationalizes it, and the case study shows how teams used
  it in practice.\par
  \medskip
  Speaker notes: This is the single slide to anchor the whole talk. Walk the
  ArchHypo framework: an engineer identifies architectural hypotheses from
  requirements and solutions; each hypothesis influences architectural
  decisions that evolve the software architecture; a technical plan manages
  the impact and uncertainty of each hypothesis in an iterative cycle. Then
  point to the right: Hypo-Stage is the plugin that supports this technique,
  and the whole study asks how engineers use Hypo-Stage to apply ArchHypo in
  a real-world organization. Everything that follows zooms into one part of
  this picture.%
}


% ============================================================================
% Slide 3 — Problem & Research Questions
% ============================================================================
% SCRIPT: "The open problem is not uncertainty itself, but the lack of
%          tool-supported, observed real-world practice --- the gap that RQ1
%          through RQ4 set out to close."
\begin{frame}{Problem \& Research Questions}
  \begin{columns}[T]
    \begin{column}{0.46\textwidth}
      \begin{block}{The Problem}
        \scriptsize
        \begin{itemize}\setlength\itemsep{0.45em}
          \item Deriving architectures still depends on architects'
                \alert{tacit knowledge} and lacks \alert{tool support}
                \citep{souza2019}.
          \item ArchHypo structures uncertainty, but teams perceive it as \alert{hard to learn}
                and missed guidance via supporting tools \citep{silva2024tse,silva2024ieeesw}.
          \item Hypo-Stage gives the initial tool support, but 
                \alert{impact on real teams} and refinements remained
                \alert{unknown} \citep{correa2025}.
        \end{itemize}
      \end{block}
    \end{column}
    \begin{column}{0.52\textwidth}
      \begin{block}{Research questions}
        \footnotesize
        \textbf{RQ1.} How do teams \emph{use} Hypo-Stage and ArchHypo on a
        scaled product?\\[0.6em]
        \textbf{RQ2.} What \emph{uncertainties} and recurring \emph{patterns}
        do they record?\\[0.6em]
        \textbf{RQ3.} What \emph{benefits and challenges} do team members
        perceive?\\[0.6em]
        \textbf{RQ4.} Which \emph{refinements} to the tool emerge, and how do
        they address the challenges?
      \end{block}
    \end{column}
  \end{columns}
\end{frame}
\note{%
  Script: The problem is not only uncertainty itself, but the lack of
  practical support for teams to share and act on architectural decisions.\par
  \medskip
  Speaker notes: State the gap in one sentence: prior ArchHypo studies showed
  conceptual value but a steep learning curve and no tool support; existing
  derivation methods (Souza et al., 2019) depend on architects' tacit
  knowledge. Then read the four RQs quickly --- use, uncertainties/patterns,
  perceived value, and refinements. Note they map one-to-one onto the boxes of
  the previous diagram.%
}


% ============================================================================
% Slide 4 — Study Setup
% ============================================================================
% SCRIPT: "We studied two teams in one organization across five sprints,
%          triangulating participant observation, the tool's own hypothesis
%          records, and active-listening sessions."
\begin{frame}{Study Setup}
  \centering
  \begin{tikzpicture}[font=\footnotesize,>=Stealth]
    % --- left column: term describing each row of boxes ---
    \node[anchor=east,align=right,font=\footnotesize\bfseries] at (-6.5,1.3) {Environment};
    \node[anchor=east,align=right,font=\footnotesize\bfseries] at (-6.5,0) {Teams};
    \node[anchor=east,align=right] at (-6.5,-1.9)
      {\textbf{Timeline}\\ \scriptsize 5 sprints\\ ($\approx$9 weeks)};
    \node[anchor=east,align=right] at (-6.5,-3.65)
      {\textbf{Data sources}\\ \scriptsize triangulated};

    % --- Row Z: environment ---
    \node[draw,rounded corners,fill=black!8,minimum width=3.0cm,
          minimum height=0.8cm,align=center] at (-3.4,1.3) {Single organization};
    \node[draw,rounded corners,fill=black!8,minimum width=3.0cm,
          minimum height=0.8cm,align=center] at (0,1.3) {Scaled product};
    \node[draw,rounded corners,fill=black!8,minimum width=3.0cm,
          minimum height=0.8cm,align=center] at (3.4,1.3) {Backstage IDP};

    % --- Row A: the two teams (cases) ---
    \node[draw,rounded corners,fill=structure.fg!10,minimum width=3.2cm,
          minimum height=1cm,align=center] (prod) at (-2.1,0)
      {\textbf{Product team}\\12 engineers};
    \node[draw,rounded corners,fill=structure.fg!10,minimum width=3.2cm,
          minimum height=1cm,align=center] (plat) at (2.1,0)
      {\textbf{Platform team}\\4 engineers};

    % --- Row B: five-sprint timeline ---
    \foreach \i [evaluate=\i as \x using {-4.4+(\i-1)*2.2}] in {1,2,3,4,5}{
      \node[draw,rounded corners,fill=DarkSlateGray!12,minimum width=1.6cm,
            minimum height=0.8cm] (s\i) at (\x,-1.9) {Sprint \i};
    }
    \foreach \i/\j in {1/2,2/3,3/4,4/5}{
      \draw[->,thick] (s\i.east) -- (s\j.west);
    }

    % --- Row C: data sources arranged as a triangulation ---
    \node[draw,rounded corners,fill=SteelBlue!15,align=center,minimum width=3.2cm,
          minimum height=0.9cm] (d2) at (0,-3.0) {Artifact analysis\\(13 hypotheses)};
    \node[draw,rounded corners,fill=SteelBlue!15,align=center,minimum width=3.2cm,
          minimum height=0.9cm] (d1) at (-2.6,-4.3) {Participant\\observation};
    \node[draw,rounded corners,fill=SteelBlue!15,align=center,minimum width=3.2cm,
          minimum height=0.9cm] (d3) at (2.6,-4.3) {Active-listening\\sessions};
    \begin{scope}[on background layer]
      \draw[densely dashed,gray!70,line width=0.7pt] (d1.north) -- (d2.south);
      \draw[densely dashed,gray!70,line width=0.7pt] (d3.north) -- (d2.south);
      \draw[densely dashed,gray!70,line width=0.7pt] (d1.east)  -- (d3.west);
    \end{scope}
  \end{tikzpicture}
\end{frame}
\note{%
  Script: We observed two teams over five sprints using participant
  observation, artifacts, and active listening to understand real tool use.\par
  \medskip
  Speaker notes: Flexible multiple-case study with formative-evaluation and
  action-research elements (Robson, 2016). Setting: one organization with a
  mature Backstage deployment and a scaled product in Latin America. Two
  contrasting cases --- a 12-engineer Product team and a 4-engineer Platform
  team --- observed over five two-week sprints (~9 calendar weeks). Three data
  sources triangulate: participant observation, artifact analysis of the
  thirteen registered hypotheses and their technical plans, and open-ended
  active-listening sessions with senior technical leaders.%
}


% ============================================================================
% Slide 5 — Research method / conceptual framework
% ============================================================================
% SCRIPT: "Methodologically this is a flexible, real-world case study: context
%          shapes how the tool is used, which yields observations and artifacts
%          that in turn drive the tool's refinements."
\begin{frame}[t]{Research Method}
  \begin{center}
    \small Grounded in the \alert{flexible design} approach to real-world
    research \citep{robson2016}.
  \end{center}
  \vspace{0.2em}
  \begin{center}
    \begin{tikzpicture}
      \node[rounded corners,fill=white,draw=black!25,line width=0.5pt,inner sep=7pt,
            blur shadow={shadow blur steps=6,shadow xshift=0.7pt,shadow yshift=-0.9pt}]
        {\includegraphics[height=0.64\textheight]{conceptual-framework}};
    \end{tikzpicture}
  \end{center}
\end{frame}
\note{%
  Script: The framework ties the study constructs to the research questions ---
  context shapes tool usage, which produces artifacts and observations that
  drive refinements.\par
  \medskip
  Speaker notes: This conceptual framework organizes the study's constructs and
  maps them onto the four research questions. Team and organization context
  (RQ1) shape and enable Arch-Stage tool usage; usage produces artifacts (RQ2)
  and experiment observations (RQ3); together they feed an iterative refinement
  cycle that generates tool refinements (RQ4). It is the lens through which the
  next slides report the results.%
}


% ============================================================================
% Slide 6 — Advances (1): how teams used Hypo-Stage   [RQ1]
% ============================================================================
% SCRIPT: "In practice, adoption was senior-led and event-anchored ---
%          registration concentrated in technical leaders and clustered around
%          protected sessions, with Backstage integration lowering the cost of
%          reaching the tool."
\begin{frame}{How teams used Hypo-Stage \hfill {\small\normalfont RQ1}}
  \begin{columns}[T]
    \begin{column}{0.56\textwidth}
      \includegraphics[width=\textwidth]{prioritization-dashboard}
    \end{column}
    \begin{column}{0.42\textwidth}
      \footnotesize
      \textbf{13 hypotheses}, two teams --- adoption was \alert{senior-led}
      and \alert{event-anchored}:
      \smallskip
      \begin{itemize}\setlength\itemsep{0.4em}
        \item \textbf{12 of 13} authored by senior/staff engineers
        \item Seniors used it for \alert{decision support}; junior/mid uncertainties resolved within a sprint, so went unregistered
        \item Registration clustered around protected technical sessions
        \item Backstage integration cut \alert{context-switching} cost
      \end{itemize}
    \end{column}
  \end{columns}
\end{frame}
\note{%
  Script: In practice, the tool was mainly adopted through senior-led
  interactions, which shaped how the teams registered and discussed
  hypotheses.\par
  \medskip
  Speaker notes: This answers RQ1. Hypothesis registration did not spread
  evenly across sprint work; it concentrated in senior technical leaders and
  clustered around discrete events (protected discussion time, training, the
  Platform team's Kafka failure-handling session) --- what the thesis calls
  "event-anchored" use. The dashboard shows the resulting backlog (13
  hypotheses, validated/open counts). Role-based engagement explains the skew:
  senior technical leaders treated Hypo-Stage as a decision-support instrument
  and referenced the backlog in architectural discussions, whereas mid- and
  junior engineers operate at a scope where uncertainties resolve within a
  sprint or less, leaving little perceived need to register --- not
  disengagement. Key enabler: delivering ArchHypo inside Backstage, where teams
  already work, lowered context-switching cost --- named explicitly by
  participants as a facilitating factor.%
}


% ============================================================================
% Slide 7 — Advances (2): uncertainties & emerging patterns   [RQ2]
% ============================================================================
% SCRIPT: "The thirteen registered uncertainties concentrate on compliance and
%          shared-infrastructure concerns; from them we abstracted five
%          recurring hypothesis patterns."
\begin{frame}{Architectural Uncertainties \& Patterns \hfill {\small\normalfont RQ2}}
  \begin{columns}[T]
    \begin{column}{0.44\textwidth}
      \includegraphics[width=\textwidth]{entity-hypotheses-tab}\\[0.3em]
      \scriptsize \textbf{13 hypotheses} (6 Product, 7 Platform), surfaced per
      service in the portal: \alert{8 requirements-} / \alert{5 solution-}sourced;
      \textbf{9 of 13} carry high impact. Dominant quality attributes:
      \textbf{Reliability} \& \textbf{Auditability}.
    \end{column}
    \begin{column}{0.54\textwidth}
      \footnotesize
      \textbf{Five emerging hypothesis patterns}\\
      {\scriptsize First empirically grounded answer to \emph{what shapes}
      hypotheses take in practice.}
      \smallskip
      \begin{itemize}\setlength\itemsep{0.1em}
        \item \textbf{Hypothesis Cluster} --- one platform concern $\rightarrow$ linked hypotheses
        \item \textbf{Regulatory Trigger} --- compliance obligation as tracked uncertainty
        \item \textbf{Context-Agnostic Restructuring} --- whether/how to decouple boundaries
        \item \textbf{Load-Validated} --- proven under load, implementation still open
        \item \textbf{Shared Resource Governance} --- ownership rules for shared infrastructure
      \end{itemize}
    \end{column}
  \end{columns}
\end{frame}
\note{%
  Script: The empirical data shows recurring shapes of architectural
  uncertainty, which we abstracted into new patterns.\par
  \medskip
  Speaker notes: RQ2. Thirteen hypotheses (six Product, seven Platform); eight
  were requirements-sourced (an unmet quality goal or external obligation) and
  five solution-sourced (uncertain consequences of a candidate decision). Nine
  of the thirteen carry High or Very High impact --- teams registered
  consequential concerns, not peripheral ones --- and a Very-Low-uncertainty /
  Very-High-impact cluster (H06-H08, H13) landed in the dashboard's "Can
  Postpone" segment. They concentrated on
  compliance-driven and shared-infrastructure concerns; Reliability and
  Auditability dominated the quality-attribute distribution, reflecting the
  regulated study site. From thematic analysis we named five emerging
  structural patterns (Group B). They complement the existing ArchHypo pattern
  language: prior patterns described HOW hypotheses are assessed/deferred;
  these describe WHAT KIND of hypothesis is being assessed. Stress this is
  nominative --- a vocabulary for recognition, not a frequency claim, given the
  dataset size.%
}


% ============================================================================
% Slide 8 — Advances (3): Perceived Benefits & Challenges   [RQ3]
% ============================================================================
% SCRIPT: "Beyond how the tool was used, participants converged on three clear
%          benefits and three recurring challenges of working this way."
\begin{frame}{Perceived Benefits \& Challenges \hfill {\small\normalfont RQ3}}
  \begin{columns}[T]
    \begin{column}{0.49\textwidth}
      \begin{block}{Benefits}
        \footnotesize
        \begin{itemize}\setlength\itemsep{0.55em}
          \item \textbf{Externalizes tacit knowledge} --- a written hypothesis becomes shareable beyond the original discussion
          \item \textbf{Supports prioritization} --- \alert{``Can Postpone''} justifies \emph{deliberate} deferral, not neglect
          \item \textbf{Collective memory} --- records open questions \emph{before} a decision, where ADRs cannot
        \end{itemize}
      \end{block}
    \end{column}
    \begin{column}{0.49\textwidth}
      \begin{block}{Challenges}
        \footnotesize
        \begin{itemize}\setlength\itemsep{0.55em}
          \item \textbf{Falsifiable phrasing} --- the \alert{negative formulation} is non-intuitive and easy to skip under pressure
          \item \textbf{Registration time cost} --- short-lived uncertainties go \alert{unregistered}
          \item \textbf{No team ceremonies} --- usage stayed \alert{individual}, never a shared review ritual
        \end{itemize}
      \end{block}
    \end{column}
  \end{columns}
\end{frame}
\note{%
  Script: Beyond how the tool was used, the two technical leaders from each team
  converged on three clear benefits and three recurring challenges.\par
  \medskip
  Speaker notes: This answers RQ3, drawn from the active-listening sessions with
  all four interviewed participants. Benefits: (1) the act of registering a
  hypothesis externalizes tacit architectural knowledge --- a clear statement
  with uncertainty/impact ratings makes the concern shareable with colleagues who
  were not in the room; (2) the "Where to Focus" panel, and especially the "Can
  Postpone" segment, gave them a principled vocabulary to explain deliberate
  deferral rather than neglect; (3) the backlog acts as a collective memory that
  captures open questions BEFORE a decision exists --- complementing ADRs, which
  only record decisions after the fact. Challenges: (1) the falsifiable,
  negative-formulation convention is non-intuitive and easy to skip under time
  pressure --- replicating the learning curve from prior ArchHypo evaluations;
  (2) the focused time to register a hypothesis means short-lifecycle
  uncertainties get under-represented; (3) no team-visible ceremony emerged, so
  usage remained largely individual within the observation window.%
}


% ============================================================================
% Slide 9 — Advances (4): Tool Refinements   [RQ4]
% ============================================================================
% SCRIPT: "Every refinement came from an observed friction point, so each of
%          the nine is traceable to concrete use during the study."
\begin{frame}{Tool Refinements \hfill {\small\normalfont RQ4}}
  \centering\small
  \textbf{9 refinement categories} --- including prioritization, the evolution
  chart, and form usability --- each traced to an observed friction point.
  \vspace{0.4em}
  \begin{columns}[c]
    \begin{column}{0.58\textwidth}
      \includegraphics[width=\textwidth]{evolution-chart}\\[0.2em]
      \scriptsize Evolution chart: uncertainty \& impact tracked across
      successive assessments.
    \end{column}
    \begin{column}{0.40\textwidth}
      \includegraphics[width=\textwidth]{hypothesis-form}\\[0.2em]
      \scriptsize Inline formulation guidance in the creation form lowers the
      falsifiability barrier.
    \end{column}
  \end{columns}
\end{frame}
\note{%
  Script: These refinements came directly from observed friction points, so
  each change is traceable to concrete use during the study.\par
  \medskip
  Speaker notes: RQ4. Nine refinement categories emerged during the
  observation window (not from upfront design), and several were redeployed
  and re-observed in the same study --- a worked instance of formative
  evaluation. The most consequential is the evolution chart (left): it closes
  the assessment-history gap of the initial Hypo-Stage version and lets teams
  judge whether technical-plan actions had the intended effect. The form
  (right) adds inline guidance and a prioritization dashboard to reduce the
  cognitive cost of writing falsifiable statements.%
}


% ============================================================================
% Slide 10 — Implications / Lessons Learned   [synthesis]
% ============================================================================
% SCRIPT: "The headline lesson: a tool removes friction but cannot, by itself,
%          supply the organizational conditions for broad participation."
\begin{frame}{Implications \& Lessons Learned}
  \vfill
  \begin{columns}[T]
    \begin{column}{0.31\textwidth}
      \begin{block}{Finding}
        Even with tool support, hypothesis work stayed \alert{senior-anchored}
        (12 of 13).
      \end{block}
    \end{column}
    \begin{column}{0.31\textwidth}
      \begin{block}{Implication}
        A tool removes \alert{friction} but cannot supply the
        \alert{organizational conditions} for broad participation.
      \end{block}
    \end{column}
    \begin{column}{0.31\textwidth}
      \begin{block}{Lesson}
        Ship hypothesis support \alert{inside} the developer portal ---
        \emph{extend} the workflow, don't add a new one.
      \end{block}
    \end{column}
  \end{columns}
  \vfill
  \centering\small
  \textbf{Takeaway:} distributed architectural work needs lightweight structure
  that preserves team autonomy while protecting shared decision quality.
  \vfill
\end{frame}
\note{%
  Script: The broader lesson is that distributed architectural work needs
  lightweight structures that preserve team autonomy while maintaining shared
  decision quality.\par
  \medskip
  Speaker notes: RQ3 framed as finding -> implication -> lesson. For
  practitioners: embedding Hypo-Stage in Backstage removed adoption friction
  (consistent with Team Topologies: tooling should follow team cognitive
  pathways). But tool availability is necessary, not sufficient --- broad
  participation also needs protected architectural time, statement-level
  scaffolding, and a recurring backlog-review ceremony (the Group A
  organizational patterns). For researchers: formative evaluation plus
  artifact analysis of the tool's own data was a productive combination for
  workflow-coupled artifacts.%
}


% ============================================================================
% Slide 11 — Future work / Conclusion
% ============================================================================
% SCRIPT: "The bounded claim is that integration and traceability worked; what
%          remains is broadening participation and validating these patterns
%          across other teams and settings."
\begin{frame}{Future Work \& Conclusion}
  \centering\small Four strategic directions:
  \begin{center}
  \begin{tikzpicture}[font=\small,>=Stealth,node distance=6mm]
    \node[draw,rounded corners,fill=structure.fg!10,align=center,
          minimum height=1cm,text width=2.6cm] (a) {Replicate\\across orgs};
    \node[draw,rounded corners,fill=structure.fg!10,align=center,
          minimum height=1cm,text width=2.6cm,right=of a] (b)
      {Longitudinal\\lifecycle study};
    \node[draw,rounded corners,fill=structure.fg!10,align=center,
          minimum height=1cm,text width=2.6cm,right=of b] (c)
      {Integrate with\\ADR tools};
    \node[draw,rounded corners,fill=structure.fg!10,align=center,
          minimum height=1cm,text width=2.6cm,right=of c] (d)
      {Broaden\\participation};
  \end{tikzpicture}
  \end{center}
  \vfill
  \begin{block}{Takeaway}
    Tool-mediated ArchHypo succeeds at \alert{integration and traceability};
    the remaining challenge is \alert{organizational and pedagogical}.
  \end{block}
  \vfill
  \centering\footnotesize
  \textbf{Limitations:} single organization \,$\cdot$\, 13 hypotheses
  \,$\cdot$\, $\approx$9 weeks \,$\cdot$\, researcher as participant.\\[0.3em]
  \emph{Emerging patterns, not consolidated findings.}
\end{frame}
\note{%
  Script: The next step is to refine the tool and validate whether these
  patterns and refinements generalize to other teams and settings.\par
  \medskip
  Speaker notes: State the limitations honestly first --- single organization
  with a mature Backstage, thirteen hypotheses over ~9 weeks, eleven still
  Open at study end, and the researcher's dual role. Then the roadmap:
  (1) replicate in other organizations/domains; (2) a longitudinal design to
  observe hypothesis closure and whether senior dependency relaxes over time;
  (3) integrate Hypo-Stage with ADR tools to connect pre-decision uncertainty
  to post-decision record; (4) intervention studies on broadening
  participation. Close on the bounded claim: integration and traceability
  worked; the deeper challenge is organizational and pedagogical. Thank the
  committee and invite questions.%
}


% ============================================================================
% Slide 12 — Closing / Thank you   (mirrors the title slide)
% ============================================================================
% SCRIPT: "Thank you --- I'd be glad to take your questions."
\makeatletter
{
  \definecolor{bggray}{RGB}{238,238,236}
  \setbeamercolor{background canvas}{bg=bggray}
  \setbeamercolor{background}{bg=}
  \setbeamertemplate{background}{%
    \ifcsvoid{@IMEtitleBgImage}{}{\@IMEtitleBgImage\hspace{-\paperwidth}}%
    \ifcsvoid{@IMElogos}{}{\@IMElogos}%
  }
  \begin{frame}[plain]
    \vspace*{1.5\baselineskip}
    \begin{minipage}{0.66\paperwidth}
      \raggedright
      {\Huge\bfseries\usebeamercolor[fg]{normal text}\color{fg} Thank you!}\par
      \vspace{0.6\baselineskip}
      {\large Questions \& discussion are welcome.}\par
      \vspace{1.2\baselineskip}
      {\small
        \textbf{José Gonçalves Lima Neto}\\[0.25em]
        Advisor: Prof.\ Dr.\ Paulo Roberto Miranda Meirelles\\
        Graduate Program in Computer Science --- IME-USP\\[0.5em]
        Master's Dissertation Defense \,$\cdot$\, 11 June 2026
      }
    \end{minipage}
  \end{frame}
}
\makeatother

\end{document}
